\subsection*{2.3 Geometric Interpretation}

\begin{conceptbox}
    Partial derivatives represent slopes of tangent lines to curves formed by slicing the surface:
    \begin{equation*}
        z = f(x,y)
    \end{equation*}
\end{conceptbox}

\begin{examplebox}
    For $z=x^2+y^2$, compute slopes:
    \begin{align*}
        \frac{\partial z}{\partial x} = 2x, \quad \frac{\partial z}{\partial y}=2y
    \end{align*}
    At $(1,1)$:
    \begin{align*}
        \frac{\partial z}{\partial x}=2, \quad \frac{\partial z}{\partial y}=2
    \end{align*}
    Interpretation: slope is steeper as distance from origin increases.
\end{examplebox}